\DAsection[日志、排错、备份、更新与恢复]{运维、排错、回滚与课堂实战}

\begin{frame}{9.1 线上故障按层定位}{从用户入口一路追到数据层}
  \centering
  \begin{tikzpicture}[node distance=2mm]
    \node[DAflow,minimum width=18mm] (dns) {DNS};
    \node[DAflow,minimum width=18mm,right=of dns] (tls) {TLS};
    \node[DAflow,minimum width=18mm,right=of tls] (nginx) {Nginx};
    \node[DAflow,minimum width=18mm,right=of nginx] (front) {Frontend};
    \node[DAflow,minimum width=18mm,right=of front] (api) {API};
    \node[DAflow,minimum width=18mm,right=of api] (db) {Database};
    \draw[DAarrowred] (dns) -- (tls);
    \draw[DAarrowred] (tls) -- (nginx);
    \draw[DAarrowred] (nginx) -- (front);
    \draw[DAarrowred] (front) -- (api);
    \draw[DAarrowred] (api) -- (db);
  \end{tikzpicture}
  \par\medskip
  \begin{itemize}
    \item 每层都用独立工具验证，找到“第一次失败”的位置。
    \item 页面报 500 时，从首次失败的层级开始处理。
  \end{itemize}
\end{frame}

\begin{frame}[fragile]{9.2 502 Bad Gateway 决策树}{Nginx 到上游的连接失败}
  \begin{lstlisting}[language=bash]
# 1. Is the process/container running?
docker compose ps
docker compose logs --tail 200 api

# 2. Is the port listening?
ss -lntp | grep 8000

# 3. Can the host/proxy reach it?
curl -i http://127.0.0.1:8000/health

# 4. Is Nginx config valid?
nginx -t
  \end{lstlisting}
  \DAtagline{结论}{本机 curl 都失败，先修应用；本机成功但域名失败，再查 Nginx 与防火墙}
\end{frame}

\begin{frame}{9.3 数据库连接失败怎么分层}{网络、账号、Schema、迁移}
  \begin{enumerate}
    \item 主机名是否正确：容器内用 \DAfile{db}，宿主机进程可能用 \DAfile{127.0.0.1}。
    \item 端口是否正确，数据库是否真正 Ready。
    \item 用户、密码、数据库名是否与初始化配置一致。
    \item 用户是否有目标数据库权限，字符集是否符合项目要求。
    \item 当前代码需要的迁移是否已执行。
  \end{enumerate}
  \DAtagline{容器陷阱}{API 容器里的 \DAfile{localhost:3306} 指向 API 容器自身}
\end{frame}

\begin{frame}[fragile]{9.4 服务器资源不足的症状}{构建时和运行时的瓶颈不同}
  \begin{lstlisting}[language=bash]
free -h
df -h
du -sh /var/lib/docker/*
docker stats
top
  \end{lstlisting}
  \begin{itemize}
    \item Node 前端构建常见内存峰值；MySQL、模型服务与多 Worker 会长期占内存。
    \item 磁盘可能被 Docker Image、Build Cache、日志和数据库共同吃满。
    \item OOM Kill 时应用像“突然消失”，需要查系统日志和容器退出码。
  \end{itemize}
\end{frame}

\begin{frame}{9.5 单机部署的安全基线}{默认拒绝，按需开放}
  \begin{columns}[T]
    \begin{column}{0.48\textwidth}
      \begin{itemize}
        \item SSH 使用密钥，限制 root 登录
        \item 面板入口限制来源并启用 HTTPS
        \item 公网只开放 22 / 80 / 443
        \item MySQL、Redis、API 端口不暴露公网
        \item 运行进程使用非 root 用户
      \end{itemize}
    \end{column}
    \begin{column}{0.48\textwidth}
      \begin{itemize}
        \item 系统与依赖定期更新
        \item Secret 最小权限并定期轮换
        \item 数据库与上传文件定期备份
        \item 日志与磁盘空间设告警
        \item 删除离职成员与过期 Deploy Key
      \end{itemize}
    \end{column}
  \end{columns}
  \DAtagline{现实边界}{一台服务器适合小型项目和教学；业务增长后再拆数据库、对象存储与负载均衡}
\end{frame}

\begin{frame}{9.6 一次规范发布包含哪些阶段}{发布前、中、后都要有证据}
  \begin{DAtable}{@{}p{23mm}p{94mm}@{}}
    \toprule
    阶段 & 动作 \\
    \midrule
    准备 & 分支已合并、测试通过、记录版本、备份、确认回滚点 \\
    构建 & 锁文件安装、生产构建、镜像打 Tag、扫描敏感信息 \\
    迁移 & 执行兼容性数据库迁移，失败则停止 \\
    发布 & 重启 / 替换进程或容器，保持变更范围最小 \\
    验证 & Health、首页、登录、数据库、关键 API、日志 \\
    观察 & 错误率、响应时间、资源、用户反馈 \\
    收尾 & 记录结果，删除临时权限，更新变更日志 \\
    \bottomrule
  \end{DAtable}
\end{frame}

\begin{frame}{9.7 回滚触发条件}{提前定义阈值与决策人}
  \begin{itemize}
    \item 健康检查持续失败，服务无法启动。
    \item 登录、支付、核心对话等关键链路不可用。
    \item 错误率或响应时间明显超过基线。
    \item 数据写入异常、权限越权、Secret 泄漏等安全事件。
    \item 修复时间不可控，继续在线排查会扩大影响。
  \end{itemize}
  \DAtagline{决策原则}{先恢复服务，再分析根因；回滚本身也需要验证与记录}
\end{frame}

\begin{frame}{9.8 从手工部署走向 CI/CD}{今天先把手工 SOP 做对，再自动化}
  \centering
  \begin{tikzpicture}[node distance=2mm]
    \node[DAflow,minimum width=18mm] (push) {Push / 合并};
    \node[DAflow,minimum width=18mm,right=of push] (ci) {CI\\Lint + Test};
    \node[DAflow,minimum width=18mm,right=of ci] (build) {Build\\Image + SHA};
    \node[DAflow,minimum width=18mm,right=of build] (registry) {Registry\\镜像};
    \node[DAflow,minimum width=18mm,right=of registry] (deploy) {Deploy\\Pull + Up};
    \node[DAflow,minimum width=18mm,right=of deploy] (verify) {Verify\\Smoke};
    \draw[DAarrowred] (push) -- (ci);
    \draw[DAarrowred] (ci) -- (build);
    \draw[DAarrowred] (build) -- (registry);
    \draw[DAarrowred] (registry) -- (deploy);
    \draw[DAarrowred] (deploy) -- (verify);
  \end{tikzpicture}
  \par\medskip
  \begin{itemize}
    \item 自动化把已验证的 SOP 固化为可重复流程。
    \item CI 使用与本地一致的 pnpm / uv 版本与锁文件策略。
  \end{itemize}
\end{frame}

\begin{frame}{9.9 课堂综合实战}{把 Day 3 / Day 4 项目真正交付给另一个人访问}
  \begin{enumerate}
    \item 确认上线前修复已经合并到主分支。
    \item 购买服务器后，完成宝塔基础环境配置。
    \item 选择宝塔手工路线或 Docker Compose 路线部署。
    \item 配置域名、\DAfile{/api} 反向代理与 SSL。
    \item 另一组使用手机网络访问并完成关键流程测试。
    \item 主动制造一次错误端口，按日志定位并恢复。
    \item 更新一个版本，再回滚到上一个版本。
  \end{enumerate}
\end{frame}

\begin{frame}{9.10 官方参考资料 · 宝塔、网络与 HTTPS}
  \scriptsize
  \begin{itemize}
    \item \href{https://docs.bt.cn/}{宝塔面板官方文档}
    \item \href{https://docs.bt.cn/user-guide/site/php/site-config/reverse-proxy}{宝塔 · 反向代理}
    \item \href{https://docs.bt.cn/getting-started/deploy-ssl}{宝塔 · 部署 SSL 证书}
    \item \href{https://docs.bt.cn/practical-tutorials/nodejs-pm2-deployment}{宝塔 · Node.js PM2 部署}
    \item \href{https://docs.nginx.com/nginx/admin-guide/web-server/reverse-proxy/}{NGINX Docs · Reverse Proxy}
    \item \href{https://developers.cloudflare.com/dns/}{Cloudflare Docs · DNS}
    \item \href{https://letsencrypt.org/how-it-works/}{Let's Encrypt · How It Works}
    \item \href{https://developer.mozilla.org/en-US/docs/Web/Security/Defenses/Mixed_content}{MDN · Mixed Content}
    \item \href{https://www.miit.gov.cn/}{工业和信息化部 · 互联网信息服务备案相关规定}
  \end{itemize}
\end{frame}
